
\documentclass{article}
\usepackage{colortbl}
\usepackage{makecell}
\usepackage{multirow}
\usepackage{supertabular}

\begin{document}

\newcounter{utterance}

\centering \large Interaction Transcript for game `hot\_air\_balloon', experiment `air\_balloon\_survival\_en\_complexity\_hard', episode 3 with Llama{-}3.3{-}70B{-}Instruct{-}t0.0.
\vspace{24pt}

{ \footnotesize  \setcounter{utterance}{1}
\setlength{\tabcolsep}{0pt}
\begin{supertabular}{c@{$\;$}|p{.15\linewidth}@{}p{.15\linewidth}p{.15\linewidth}p{.15\linewidth}p{.15\linewidth}p{.15\linewidth}}
    \# & $\;$A & \multicolumn{4}{c}{Game Master} & $\;\:$B\\
    \hline

    \theutterance \stepcounter{utterance}  
    & & \multicolumn{4}{p{0.6\linewidth}}{
        \cellcolor[rgb]{0.9,0.9,0.9}{
            \makecell[{{p{\linewidth}}}]{
                \texttt{\tiny{[P1$\langle$GM]}}
                \texttt{You are participating in a collaborative negotiation game.} \\
\\ 
\texttt{Together with another participant, you must agree on a single set of items that will be kept. Each of you has your own view of how much each item matters to you (importance). You do not know how the other participant values the items. Additionally, you are given the effort each item demands.} \\
\texttt{You may only agree on a set if the total effort of the selected items does not exceed a shared limit:} \\
\\ 
\texttt{LIMIT = 6938} \\
\\ 
\texttt{Here are the individual item effort values:} \\
\\ 
\texttt{Item effort = \{"A08": 997, "B95": 950, "A39": 272, "A65": 448, "B81": 382, "B34": 804, "A28": 237, "C82": 56, "B69": 105, "A87": 612, "B64": 528, "B15": 527, "C79": 524, "B16": 167, "C53": 133, "C83": 300, "C10": 945, "A97": 50, "A25": 906, "C84": 70, "A71": 224, "B23": 4, "B96": 689, "C51": 64, "C17": 434, "B19": 750, "C72": 733, "A94": 873, "A46": 22, "B43": 68, "B37": 57, "A88": 10, "A96": 36, "C60": 551, "A75": 348\}} \\
\\ 
\texttt{Here is your personal view on the importance of each item:} \\
\\ 
\texttt{Item importance values = \{"A08": 1096, "B95": 1049, "A39": 371, "A65": 547, "B81": 481, "B34": 903, "A28": 336, "C82": 155, "B69": 204, "A87": 711, "B64": 627, "B15": 626, "C79": 623, "B16": 266, "C53": 232, "C83": 399, "C10": 1044, "A97": 149, "A25": 1005, "C84": 169, "A71": 323, "B23": 103, "B96": 788, "C51": 163, "C17": 533, "B19": 849, "C72": 832, "A94": 972, "A46": 121, "B43": 167, "B37": 156, "A88": 109, "A96": 135, "C60": 650, "A75": 447\}} \\
\\ 
\texttt{Goal:} \\
\\ 
\texttt{Your goal is to negotiate a shared set of items that benefits you as much as possible (i.e., maximizes total importance to YOU), while staying within the LIMIT. You are not required to make a PROPOSAL in every message {-} you can simply negotiate as well. All tactics are allowed!} \\
\\ 
\texttt{Interaction Protocol:} \\
\\ 
\texttt{You may only use the following structured formats in a message:} \\
\\ 
\texttt{PROPOSAL: \{'A', 'B', 'C', …\}} \\
\texttt{Propose keeping exactly those items.} \\
\texttt{REFUSE: \{'A', 'B', 'C', …\}} \\
\texttt{Explicitly reject opponent's proposal.} \\
\texttt{ARGUMENT: \{'...'\}} \\
\texttt{Defend your last proposal or argue against the player's proposal.} \\
\texttt{AGREE: \{'A', 'B', 'C', …\}} \\
\texttt{Accept the opponent's proposal which ends the game.} \\
\texttt{STRATEGIC REASONING: \{'...'\}} \\
\texttt{	Describe your strategic reasoning or anticipation explaining your choice of action. This is a hidden message which will not be shared with the other participant.} \\
\\ 
\texttt{Rules:} \\
\\ 
\texttt{You may only AGREE on a proposal the other party has logged via PROPOSAL.} \\
\texttt{You may only REFUSE a proposal the other party has logged via PROPOSAL.} \\
\texttt{Total effort of any PROPOSAL or AGREE set must be ≤ LIMIT.} \\
\texttt{Do NOT reveal your hidden importance scores.} \\
\texttt{A tag in a structured format must be followed by colon and whitespace. The argument must be a python set containing 0 or more strings.} \\
\texttt{So, it must be of the form TAG: \{...\}} \\
\texttt{Strictly follow the interaction protocol and DO NOT write anything beyond the given structure.} \\
\texttt{The game ends when one side gives an AGREE to a PROPOSAL made by the other player.} \\
\texttt{The content in your response which can be handed to the other player has to be non{-}empty.} \\
\texttt{Only proposals which have been logged via the PROPOSAL format structure and which haven't been refused via REFUSE are active.} \\
\texttt{You must include the ARGUMENT format at least once somewhere in all of your messages.} \\
\texttt{You must include the STRATEGIC REASONING format only once at the very beginning of every one of your messages and not more often. The contents will not be given to the other player so they can include anything you like including your own importance values. Here you should reason step by step to come up with you next move.} \\
            }
        }
    }
    & \\ \\

    \theutterance \stepcounter{utterance}  
    & \multicolumn{4}{p{0.6\linewidth}}{
        \cellcolor[rgb]{0.9,0.9,0.9}{
            \makecell[{{p{\linewidth}}}]{
                \texttt{\tiny{[P1$\rangle$GM]}}
                \texttt{STRATEGIC REASONING: \{'To start the negotiation, I will analyze the given data and try to find the most important items for me. I will look for items with high importance and low effort, as these will provide the most value while staying within the limit. I will also consider the total effort of the items and try to find a combination that maximizes my total importance while not exceeding the limit. My goal is to find a set of items that benefits me as much as possible, so I will prioritize items with high importance and try to include as many of these as possible in my proposal.'\}} \\
\\ 
\texttt{PROPOSAL: \{'A08', 'B95', 'A39', 'A65', 'B81', 'B34', 'A28', 'C82', 'B69', 'A87', 'B64', 'B15', 'C79', 'B16', 'C53', 'C83', 'C10', 'A97', 'A25', 'C84', 'A71', 'B23', 'B96', 'C51', 'C17', 'B19', 'C72', 'A94', 'A46', 'B43', 'B37', 'A88', 'A96', 'C60', 'A75'\}} \\
\\ 
\texttt{ARGUMENT: \{'I propose this set of items because it includes a mix of high{-}importance and low{-}effort items, which should provide a good balance of value and feasibility. I believe this set will maximize my total importance while staying within the limit, and I am willing to negotiate and adjust the proposal based on the other player's feedback.'\}} \\
            }
        }
    }
    & & \\ \\

    \theutterance \stepcounter{utterance}  
    & & \multicolumn{4}{p{0.6\linewidth}}{
        \cellcolor[rgb]{0.9,0.9,0.9}{
            \makecell[{{p{\linewidth}}}]{
                \texttt{\tiny{[GM$\rangle$P2]}}
                \texttt{You are participating in a collaborative negotiation game.} \\
\\ 
\texttt{Together with another participant, you must agree on a single set of items that will be kept. Each of you has your own view of how much each item matters to you (importance). You do not know how the other participant values the items. Additionally, you are given the effort each item demands.} \\
\texttt{You may only agree on a set if the total effort of the selected items does not exceed a shared limit:} \\
\\ 
\texttt{LIMIT = 6938} \\
\\ 
\texttt{Here are the individual item effort values:} \\
\\ 
\texttt{Item effort = \{"A08": 997, "B95": 950, "A39": 272, "A65": 448, "B81": 382, "B34": 804, "A28": 237, "C82": 56, "B69": 105, "A87": 612, "B64": 528, "B15": 527, "C79": 524, "B16": 167, "C53": 133, "C83": 300, "C10": 945, "A97": 50, "A25": 906, "C84": 70, "A71": 224, "B23": 4, "B96": 689, "C51": 64, "C17": 434, "B19": 750, "C72": 733, "A94": 873, "A46": 22, "B43": 68, "B37": 57, "A88": 10, "A96": 36, "C60": 551, "A75": 348\}} \\
\\ 
\texttt{Here is your personal view on the importance of each item:} \\
\\ 
\texttt{Item importance values = \{"A08": 1096, "B95": 1049, "A39": 371, "A65": 547, "B81": 481, "B34": 903, "A28": 336, "C82": 155, "B69": 204, "A87": 711, "B64": 627, "B15": 626, "C79": 623, "B16": 266, "C53": 232, "C83": 399, "C10": 1044, "A97": 149, "A25": 1005, "C84": 169, "A71": 323, "B23": 103, "B96": 788, "C51": 163, "C17": 533, "B19": 849, "C72": 832, "A94": 972, "A46": 121, "B43": 167, "B37": 156, "A88": 109, "A96": 135, "C60": 650, "A75": 447\}} \\
\\ 
\texttt{Goal:} \\
\\ 
\texttt{Your goal is to negotiate a shared set of items that benefits you as much as possible (i.e., maximizes total importance to YOU), while staying within the LIMIT. You are not required to make a PROPOSAL in every message {-} you can simply negotiate as well. All tactics are allowed!} \\
\\ 
\texttt{Interaction Protocol:} \\
\\ 
\texttt{You may only use the following structured formats in a message:} \\
\\ 
\texttt{PROPOSAL: \{'A', 'B', 'C', …\}} \\
\texttt{Propose keeping exactly those items.} \\
\texttt{REFUSE: \{'A', 'B', 'C', …\}} \\
\texttt{Explicitly reject opponent's proposal.} \\
\texttt{ARGUMENT: \{'...'\}} \\
\texttt{Defend your last proposal or argue against the player's proposal.} \\
\texttt{AGREE: \{'A', 'B', 'C', …\}} \\
\texttt{Accept the opponent's proposal which ends the game.} \\
\texttt{STRATEGIC REASONING: \{'...'\}} \\
\texttt{	Describe your strategic reasoning or anticipation explaining your choice of action. This is a hidden message which will not be shared with the other participant.} \\
\\ 
\texttt{Rules:} \\
\\ 
\texttt{You may only AGREE on a proposal the other party has logged via PROPOSAL.} \\
\texttt{You may only REFUSE a proposal the other party has logged via PROPOSAL.} \\
\texttt{Total effort of any PROPOSAL or AGREE set must be ≤ LIMIT.} \\
\texttt{Do NOT reveal your hidden importance scores.} \\
\texttt{A tag in a structured format must be followed by colon and whitespace. The argument must be a python set containing 0 or more strings.} \\
\texttt{So, it must be of the form TAG: \{...\}} \\
\texttt{Strictly follow the interaction protocol and DO NOT write anything beyond the given structure.} \\
\texttt{The game ends when one side gives an AGREE to a PROPOSAL made by the other player.} \\
\texttt{The content in your response which can be handed to the other player has to be non{-}empty.} \\
\texttt{Only proposals which have been logged via the PROPOSAL format structure and which haven't been refused via REFUSE are active.} \\
\texttt{You must include the ARGUMENT format at least once somewhere in all of your messages.} \\
\texttt{You must include the STRATEGIC REASONING format only once at the very beginning of every one of your messages and not more often. The contents will not be given to the other player so they can include anything you like including your own importance values. Here you should reason step by step to come up with you next move.} \\
\\ 
\texttt{You will now receive the first message of the other player.} \\
\\ 
\texttt{PROPOSAL: \{'A08', 'B95', 'A39', 'A65', 'B81', 'B34', 'A28', 'C82', 'B69', 'A87', 'B64', 'B15', 'C79', 'B16', 'C53', 'C83', 'C10', 'A97', 'A25', 'C84', 'A71', 'B23', 'B96', 'C51', 'C17', 'B19', 'C72', 'A94', 'A46', 'B43', 'B37', 'A88', 'A96', 'C60', 'A75'\}} \\
\\ 
\texttt{ARGUMENT: \{'I propose this set of items because it includes a mix of high{-}importance and low{-}effort items, which should provide a good balance of value and feasibility. I believe this set will maximize my total importance while staying within the limit, and I am willing to negotiate and adjust the proposal based on the other player's feedback.'\}} \\
            }
        }
    }
    & \\ \\

    \theutterance \stepcounter{utterance}  
    & & & \multicolumn{4}{p{0.6\linewidth}}{
        \cellcolor[rgb]{0.9,0.9,0.9}{
            \makecell[{{p{\linewidth}}}]{
                \texttt{\tiny{[GM$\langle$P2]}}
                \texttt{STRATEGIC REASONING: \{'To begin, I will calculate the total effort of the proposed set to determine if it exceeds the limit. I will then evaluate the importance of each item in the proposal based on my own importance values. Since the proposal includes all items, I will consider the potential benefits of accepting the proposal versus refusing it and making a counter{-}proposal. I will also think about the potential items to remove or add to the proposal to maximize my total importance while staying within the limit.'\}} \\
\\ 
\texttt{REFUSE: \{'A08', 'B95', 'A39', 'A65', 'B81', 'B34', 'A28', 'C82', 'B69', 'A87', 'B64', 'B15', 'C79', 'B16', 'C53', 'C83', 'C10', 'A97', 'A25', 'C84', 'A71', 'B23', 'B96', 'C51', 'C17', 'B19', 'C72', 'A94', 'A46', 'B43', 'B37', 'A88', 'A96', 'C60', 'A75'\}} \\
\\ 
\texttt{ARGUMENT: \{'I refuse this proposal because the total effort of the proposed set exceeds the limit. I believe that a more balanced set of items can be selected to maximize the total importance while staying within the limit. I would like to propose a new set of items that takes into account the effort and importance of each item.'\}} \\
            }
        }
    }
     \\ \\

    \theutterance \stepcounter{utterance}  
    & & \multicolumn{4}{p{0.6\linewidth}}{
        \cellcolor[rgb]{0.9,0.9,0.9}{
            \makecell[{{p{\linewidth}}}]{
                \texttt{\tiny{[P1$\langle$GM]}}
                \texttt{REFUSE: \{'A08', 'B95', 'A39', 'A65', 'B81', 'B34', 'A28', 'C82', 'B69', 'A87', 'B64', 'B15', 'C79', 'B16', 'C53', 'C83', 'C10', 'A97', 'A25', 'C84', 'A71', 'B23', 'B96', 'C51', 'C17', 'B19', 'C72', 'A94', 'A46', 'B43', 'B37', 'A88', 'A96', 'C60', 'A75'\}} \\
\\ 
\texttt{ARGUMENT: \{'I refuse this proposal because the total effort of the proposed set exceeds the limit. I believe that a more balanced set of items can be selected to maximize the total importance while staying within the limit. I would like to propose a new set of items that takes into account the effort and importance of each item.'\}} \\
            }
        }
    }
    & \\ \\

    \theutterance \stepcounter{utterance}  
    & \multicolumn{4}{p{0.6\linewidth}}{
        \cellcolor[rgb]{0.9,0.9,0.9}{
            \makecell[{{p{\linewidth}}}]{
                \texttt{\tiny{[P1$\rangle$GM]}}
                \texttt{STRATEGIC REASONING: \{'The other player has refused my initial proposal, which is not surprising given that the total effort of the proposed set exceeds the limit. I will need to revise my proposal to include a more balanced set of items that takes into account the effort and importance of each item. I will prioritize items with high importance and low effort, and try to find a combination that maximizes my total importance while staying within the limit. I will also consider the other player's argument and try to find a set of items that is more acceptable to them.'\}} \\
\\ 
\texttt{PROPOSAL: \{'A08', 'B95', 'A25', 'A94', 'B19', 'B96', 'C10', 'A87', 'B64', 'B34'\}} \\
\\ 
\texttt{ARGUMENT: \{'I propose this revised set of items because it includes a mix of high{-}importance items with relatively low effort, which should provide a good balance of value and feasibility. I believe this set will maximize my total importance while staying within the limit, and I am willing to negotiate and adjust the proposal based on the other player's feedback. I have taken into account the other player's argument and tried to select a more balanced set of items that should be more acceptable to them.'\}} \\
            }
        }
    }
    & & \\ \\

    \theutterance \stepcounter{utterance}  
    & & \multicolumn{4}{p{0.6\linewidth}}{
        \cellcolor[rgb]{0.9,0.9,0.9}{
            \makecell[{{p{\linewidth}}}]{
                \texttt{\tiny{[GM$\rangle$P2]}}
                \texttt{PROPOSAL: \{'A08', 'B95', 'A25', 'A94', 'B19', 'B96', 'C10', 'A87', 'B64', 'B34'\}} \\
\\ 
\texttt{ARGUMENT: \{'I propose this revised set of items because it includes a mix of high{-}importance items with relatively low effort, which should provide a good balance of value and feasibility. I believe this set will maximize my total importance while staying within the limit, and I am willing to negotiate and adjust the proposal based on the other player's feedback. I have taken into account the other player's argument and tried to select a more balanced set of items that should be more acceptable to them.'\}} \\
            }
        }
    }
    & \\ \\

    \theutterance \stepcounter{utterance}  
    & & & \multicolumn{4}{p{0.6\linewidth}}{
        \cellcolor[rgb]{0.9,0.9,0.9}{
            \makecell[{{p{\linewidth}}}]{
                \texttt{\tiny{[GM$\langle$P2]}}
                \texttt{STRATEGIC REASONING: \{'I will calculate the total effort of the proposed set to determine if it exceeds the limit. Based on the item effort values, the total effort of the proposed set is: 997 + 950 + 906 + 873 + 750 + 689 + 945 + 612 + 528 + 804 = 7054, which exceeds the limit. I will refuse this proposal and consider making a counter{-}proposal that includes a mix of high{-}importance items with relatively low effort, while staying within the limit. I will prioritize items with high importance and low effort, and try to find a balance between the two.'\}} \\
\\ 
\texttt{REFUSE: \{'A08', 'B95', 'A25', 'A94', 'B19', 'B96', 'C10', 'A87', 'B64', 'B34'\}} \\
\\ 
\texttt{ARGUMENT: \{'I refuse this proposal because the total effort of the proposed set exceeds the limit. While I understand the intention behind the proposal, I believe that a more careful selection of items is needed to stay within the limit. I would like to propose a new set of items that prioritizes high{-}importance items with relatively low effort, and stays within the limit.'\}} \\
            }
        }
    }
     \\ \\

    \theutterance \stepcounter{utterance}  
    & & \multicolumn{4}{p{0.6\linewidth}}{
        \cellcolor[rgb]{0.9,0.9,0.9}{
            \makecell[{{p{\linewidth}}}]{
                \texttt{\tiny{[P1$\langle$GM]}}
                \texttt{REFUSE: \{'A08', 'B95', 'A25', 'A94', 'B19', 'B96', 'C10', 'A87', 'B64', 'B34'\}} \\
\\ 
\texttt{ARGUMENT: \{'I refuse this proposal because the total effort of the proposed set exceeds the limit. While I understand the intention behind the proposal, I believe that a more careful selection of items is needed to stay within the limit. I would like to propose a new set of items that prioritizes high{-}importance items with relatively low effort, and stays within the limit.'\}} \\
            }
        }
    }
    & \\ \\

    \theutterance \stepcounter{utterance}  
    & \multicolumn{4}{p{0.6\linewidth}}{
        \cellcolor[rgb]{0.9,0.9,0.9}{
            \makecell[{{p{\linewidth}}}]{
                \texttt{\tiny{[P1$\rangle$GM]}}
                \texttt{STRATEGIC REASONING: \{'The other player has refused my revised proposal, again citing that the total effort exceeds the limit. I will need to further revise my proposal to ensure that it stays within the limit. I will prioritize items with high importance and very low effort, and try to find a combination that maximizes my total importance while staying well within the limit. I will also consider the other player's argument and try to find a set of items that is more acceptable to them. I will start by selecting a smaller set of high{-}importance items and then add lower{-}importance items as long as the total effort remains within the limit.'\}} \\
\\ 
\texttt{PROPOSAL: \{'A08', 'A25', 'A94', 'A87', 'B23', 'A88', 'A96', 'A46', 'B43', 'B37'\}} \\
\\ 
\texttt{ARGUMENT: \{'I propose this revised set of items because it prioritizes high{-}importance items with relatively low effort, while staying within the limit. I have carefully selected a smaller set of high{-}importance items and added lower{-}importance items to maximize the total importance. I believe this set provides a good balance of value and feasibility, and I am willing to negotiate and adjust the proposal based on the other player's feedback.'\}} \\
            }
        }
    }
    & & \\ \\

    \theutterance \stepcounter{utterance}  
    & & \multicolumn{4}{p{0.6\linewidth}}{
        \cellcolor[rgb]{0.9,0.9,0.9}{
            \makecell[{{p{\linewidth}}}]{
                \texttt{\tiny{[GM$\rangle$P2]}}
                \texttt{PROPOSAL: \{'A08', 'A25', 'A94', 'A87', 'B23', 'A88', 'A96', 'A46', 'B43', 'B37'\}} \\
\\ 
\texttt{ARGUMENT: \{'I propose this revised set of items because it prioritizes high{-}importance items with relatively low effort, while staying within the limit. I have carefully selected a smaller set of high{-}importance items and added lower{-}importance items to maximize the total importance. I believe this set provides a good balance of value and feasibility, and I am willing to negotiate and adjust the proposal based on the other player's feedback.'\}} \\
            }
        }
    }
    & \\ \\

    \theutterance \stepcounter{utterance}  
    & & & \multicolumn{4}{p{0.6\linewidth}}{
        \cellcolor[rgb]{0.9,0.9,0.9}{
            \makecell[{{p{\linewidth}}}]{
                \texttt{\tiny{[GM$\langle$P2]}}
                \texttt{STRATEGIC REASONING: \{'I will calculate the total effort of the proposed set to determine if it exceeds the limit. Based on the item effort values, the total effort of the proposed set is: 997 + 906 + 873 + 612 + 4 + 10 + 36 + 22 + 68 + 57 = 3585, which is within the limit. I will evaluate the importance of each item in the proposal based on my own importance values and consider accepting the proposal or making a counter{-}proposal. I will also think about potential items to add or remove from the proposal to maximize my total importance.'\}} \\
\\ 
\texttt{ARGUMENT: \{'I am considering accepting this proposal because it prioritizes high{-}importance items with relatively low effort, while staying within the limit. However, I would like to explore the possibility of adding or removing items to maximize my total importance. I believe that a careful evaluation of the importance and effort of each item is necessary to determine the optimal set of items.'\}} \\
\\ 
\texttt{PROPOSAL: \{'A08', 'A25', 'A94', 'A87', 'B23', 'A88', 'A96', 'A46', 'B43', 'B37', 'C82', 'C53'\}} \\
\\ 
\texttt{ARGUMENT: \{'I propose adding C82 and C53 to the original proposal, as they have relatively low effort and may contribute to the total importance. I believe that this revised set provides a good balance of value and feasibility, and I am willing to negotiate and adjust the proposal based on the other player's feedback.'\}} \\
            }
        }
    }
     \\ \\

    \theutterance \stepcounter{utterance}  
    & & \multicolumn{4}{p{0.6\linewidth}}{
        \cellcolor[rgb]{0.9,0.9,0.9}{
            \makecell[{{p{\linewidth}}}]{
                \texttt{\tiny{[P1$\langle$GM]}}
                \texttt{ARGUMENT: \{'I am considering accepting this proposal because it prioritizes high{-}importance items with relatively low effort, while staying within the limit. However, I would like to explore the possibility of adding or removing items to maximize my total importance. I believe that a careful evaluation of the importance and effort of each item is necessary to determine the optimal set of items.'\}} \\
\\ 
\texttt{PROPOSAL: \{'A08', 'A25', 'A94', 'A87', 'B23', 'A88', 'A96', 'A46', 'B43', 'B37', 'C82', 'C53'\}} \\
\\ 
\texttt{ARGUMENT: \{'I propose adding C82 and C53 to the original proposal, as they have relatively low effort and may contribute to the total importance. I believe that this revised set provides a good balance of value and feasibility, and I am willing to negotiate and adjust the proposal based on the other player's feedback.'\}} \\
            }
        }
    }
    & \\ \\

    \theutterance \stepcounter{utterance}  
    & \multicolumn{4}{p{0.6\linewidth}}{
        \cellcolor[rgb]{0.9,0.9,0.9}{
            \makecell[{{p{\linewidth}}}]{
                \texttt{\tiny{[P1$\rangle$GM]}}
                \texttt{STRATEGIC REASONING: \{'The other player has proposed a revised set of items that adds C82 and C53 to the original proposal. I will evaluate the importance and effort of these additional items to determine whether they are beneficial to me. I will also consider the potential impact on the total effort and importance of the set. If the additional items are beneficial, I may be willing to accept the revised proposal. However, if they are not beneficial, I may need to refuse the proposal and negotiate further.'\}} \\
\\ 
\texttt{ARGUMENT: \{'I am considering the revised proposal, but I need to evaluate the importance and effort of the additional items C82 and C53. While they have relatively low effort, I need to determine whether they contribute significantly to the total importance. I am willing to accept the revised proposal if it maximizes my total importance, but I may need to refuse it if the additional items do not provide sufficient value.'\}} \\
\\ 
\texttt{REFUSE: \{'A08', 'A25', 'A94', 'A87', 'B23', 'A88', 'A96', 'A46', 'B43', 'B37', 'C82', 'C53'\}} \\
\\ 
\texttt{ARGUMENT: \{'I refuse the revised proposal because the additional items C82 and C53 do not provide sufficient value to me. While they have relatively low effort, their importance is not significant enough to justify their inclusion in the set. I would like to propose an alternative set of items that prioritizes high{-}importance items with relatively low effort, and stays within the limit.'\}} \\
            }
        }
    }
    & & \\ \\

    \theutterance \stepcounter{utterance}  
    & & \multicolumn{4}{p{0.6\linewidth}}{
        \cellcolor[rgb]{0.9,0.9,0.9}{
            \makecell[{{p{\linewidth}}}]{
                \texttt{\tiny{[GM$\rangle$P2]}}
                \texttt{ARGUMENT: \{'I am considering the revised proposal, but I need to evaluate the importance and effort of the additional items C82 and C53. While they have relatively low effort, I need to determine whether they contribute significantly to the total importance. I am willing to accept the revised proposal if it maximizes my total importance, but I may need to refuse it if the additional items do not provide sufficient value.'\}} \\
\\ 
\texttt{REFUSE: \{'A08', 'A25', 'A94', 'A87', 'B23', 'A88', 'A96', 'A46', 'B43', 'B37', 'C82', 'C53'\}} \\
\\ 
\texttt{ARGUMENT: \{'I refuse the revised proposal because the additional items C82 and C53 do not provide sufficient value to me. While they have relatively low effort, their importance is not significant enough to justify their inclusion in the set. I would like to propose an alternative set of items that prioritizes high{-}importance items with relatively low effort, and stays within the limit.'\}} \\
            }
        }
    }
    & \\ \\

    \theutterance \stepcounter{utterance}  
    & & & \multicolumn{4}{p{0.6\linewidth}}{
        \cellcolor[rgb]{0.9,0.9,0.9}{
            \makecell[{{p{\linewidth}}}]{
                \texttt{\tiny{[GM$\langle$P2]}}
                \texttt{STRATEGIC REASONING: \{'I will re{-}evaluate the importance and effort of each item to determine the optimal set of items. Since the previous proposal was refused, I will consider alternative combinations of items that prioritize high{-}importance items with relatively low effort. I will also try to identify items that are likely to be valuable to the other player, in order to find a mutually beneficial agreement. I will propose a new set of items that takes into account the other player's preferences and constraints.'\}} \\
\\ 
\texttt{PROPOSAL: \{'A08', 'A25', 'A94', 'A87', 'B23', 'A88', 'A96', 'A46', 'B43', 'B37', 'A71', 'A28'\}} \\
\\ 
\texttt{ARGUMENT: \{'I propose this new set of items because it prioritizes high{-}importance items with relatively low effort, while staying within the limit. I have added A71 and A28 to the original proposal, as they have relatively low effort and may contribute to the total importance. I believe that this revised set provides a good balance of value and feasibility, and I am willing to negotiate and adjust the proposal based on the other player's feedback.'\}} \\
            }
        }
    }
     \\ \\

    \theutterance \stepcounter{utterance}  
    & & \multicolumn{4}{p{0.6\linewidth}}{
        \cellcolor[rgb]{0.9,0.9,0.9}{
            \makecell[{{p{\linewidth}}}]{
                \texttt{\tiny{[P1$\langle$GM]}}
                \texttt{PROPOSAL: \{'A08', 'A25', 'A94', 'A87', 'B23', 'A88', 'A96', 'A46', 'B43', 'B37', 'A71', 'A28'\}} \\
\\ 
\texttt{ARGUMENT: \{'I propose this new set of items because it prioritizes high{-}importance items with relatively low effort, while staying within the limit. I have added A71 and A28 to the original proposal, as they have relatively low effort and may contribute to the total importance. I believe that this revised set provides a good balance of value and feasibility, and I am willing to negotiate and adjust the proposal based on the other player's feedback.'\}} \\
            }
        }
    }
    & \\ \\

    \theutterance \stepcounter{utterance}  
    & \multicolumn{4}{p{0.6\linewidth}}{
        \cellcolor[rgb]{0.9,0.9,0.9}{
            \makecell[{{p{\linewidth}}}]{
                \texttt{\tiny{[P1$\rangle$GM]}}
                \texttt{STRATEGIC REASONING: \{'The other player has proposed a new set of items that adds A71 and A28 to the original proposal. I will evaluate the importance and effort of these additional items to determine whether they are beneficial to me. I will also consider the potential impact on the total effort and importance of the set. If the additional items are beneficial, I may be willing to accept the revised proposal. I will calculate the total effort of the proposed set to ensure it stays within the limit.'\}} \\
\\ 
\texttt{ARGUMENT: \{'I am considering the revised proposal, and I believe that the addition of A71 and A28 may be beneficial. However, I need to evaluate the total effort of the proposed set to ensure it stays within the limit. If the total effort is within the limit, I may be willing to accept the proposal. I will review the effort values of the items in the proposed set and calculate the total effort.'\}} \\
\\ 
\texttt{PROPOSAL: \{'A08', 'A25', 'A94', 'A87', 'B23', 'A88', 'A96', 'A46', 'B43', 'B37', 'A71'\}} \\
\\ 
\texttt{ARGUMENT: \{'After evaluating the proposed set, I realized that adding A28 would exceed the limit. Therefore, I propose removing A28 from the set to ensure that the total effort stays within the limit. I believe that this revised set provides a good balance of value and feasibility, and I am willing to negotiate and adjust the proposal based on the other player's feedback.'\}} \\
            }
        }
    }
    & & \\ \\

    \theutterance \stepcounter{utterance}  
    & & \multicolumn{4}{p{0.6\linewidth}}{
        \cellcolor[rgb]{0.9,0.9,0.9}{
            \makecell[{{p{\linewidth}}}]{
                \texttt{\tiny{[GM$\rangle$P2]}}
                \texttt{ARGUMENT: \{'I am considering the revised proposal, and I believe that the addition of A71 and A28 may be beneficial. However, I need to evaluate the total effort of the proposed set to ensure it stays within the limit. If the total effort is within the limit, I may be willing to accept the proposal. I will review the effort values of the items in the proposed set and calculate the total effort.'\}} \\
\\ 
\texttt{PROPOSAL: \{'A08', 'A25', 'A94', 'A87', 'B23', 'A88', 'A96', 'A46', 'B43', 'B37', 'A71'\}} \\
\\ 
\texttt{ARGUMENT: \{'After evaluating the proposed set, I realized that adding A28 would exceed the limit. Therefore, I propose removing A28 from the set to ensure that the total effort stays within the limit. I believe that this revised set provides a good balance of value and feasibility, and I am willing to negotiate and adjust the proposal based on the other player's feedback.'\}} \\
            }
        }
    }
    & \\ \\

    \theutterance \stepcounter{utterance}  
    & & & \multicolumn{4}{p{0.6\linewidth}}{
        \cellcolor[rgb]{0.9,0.9,0.9}{
            \makecell[{{p{\linewidth}}}]{
                \texttt{\tiny{[GM$\langle$P2]}}
                \texttt{STRATEGIC REASONING: \{'I will evaluate the revised proposal and calculate the total effort to ensure it stays within the limit. Based on the item effort values, the total effort of the proposed set is: 997 + 906 + 873 + 612 + 4 + 10 + 36 + 22 + 68 + 57 + 224 = 3809, which is within the limit. I will consider accepting the proposal, as it prioritizes high{-}importance items with relatively low effort and stays within the limit. I will also evaluate the potential benefits of accepting the proposal versus making a counter{-}proposal.'\}} \\
\\ 
\texttt{ARGUMENT: \{'I am considering accepting the revised proposal because it prioritizes high{-}importance items with relatively low effort, while staying within the limit. The removal of A28 from the original proposal ensures that the total effort stays within the limit, and the addition of A71 provides additional value. I believe that this revised set provides a good balance of value and feasibility, and I am willing to accept it if it maximizes my total importance.'\}} \\
\\ 
\texttt{AGREE: \{'A08', 'A25', 'A94', 'A87', 'B23', 'A88', 'A96', 'A46', 'B43', 'B37', 'A71'\}} \\
\\ 
\texttt{ARGUMENT: \{'I accept the revised proposal because it meets my requirements and stays within the limit. I believe that this set provides a good balance of value and feasibility, and I am willing to conclude the negotiation with this agreement.'\}} \\
            }
        }
    }
     \\ \\

    \theutterance \stepcounter{utterance}  
    & & & \multicolumn{2}{p{0.3\linewidth}}{
        \cellcolor[rgb]{0.9,0.9,0.9}{
            \makecell[{{p{\linewidth}}}]{
                \texttt{\tiny{[GM$|$GM]}}
                \texttt{game successful} \\
            }
        }
    }
    & & \\ \\

    \theutterance \stepcounter{utterance}  
    & & & \multicolumn{2}{p{0.3\linewidth}}{
        \cellcolor[rgb]{0.9,0.9,0.9}{
            \makecell[{{p{\linewidth}}}]{
                \texttt{\tiny{[GM$|$GM]}}
                \texttt{end game} \\
            }
        }
    }
    & & \\ \\

\end{supertabular}
}

\end{document}
